% !TeX spellcheck = pt_PT

\chapter{Conclusão}
\label{chap:8}

Este capítulo apresenta uma síntese do trabalho desenvolvido no âmbito do projeto ShopISEL e reflete sobre os principais desafios encontrados. As direções de evolução futura identificadas mas não exploradas no âmbito deste projeto, por restrições de tempo, são apresentadas em detalhe no Capítulo~\ref{chap:7}.

\section{Síntese do Trabalho Desenvolvido}
\label{sec:8.1}

O ShopISEL partiu de uma necessidade concreta dos consumidores: organizar compras num contexto de múltiplas insígnias, preços variáveis e informação dispersa. A solução desenvolvida reúne uma aplicação móvel, uma aplicação web e uma infraestrutura backend baseada em microserviços.

Em termos funcionais, a plataforma já cobre os fluxos principais do projeto: gestão de listas, consulta de preços, alertas de descida de preço, leitura de códigos de barras e apresentação de preços com base na localização do utilizador quando aplicável. Estes elementos mostram que o sistema não ficou apenas pela ideia inicial, mas evoluiu para um conjunto de funcionalidades utilizáveis em contexto real.

Do ponto de vista técnico, o maior desafio não foi apenas implementar ecrãs ou endpoints, mas ligar informação proveniente de fontes com estruturas, formatos e níveis de detalhe distintos, com autenticação centralizada, notificação push e sincronização entre serviços sem perder simplicidade de utilização.

Os principais componentes desenvolvidos foram:

\begin{itemize}
    \item Cinco serviços de backend em ASP.NET Core (listas, produtos, preços, lojas, contas), cada um com persistência relacional em PostgreSQL;
    \item Um serviço híbrido (MatchQueue) que faz a ponte entre dados de scraping armazenados em MongoDB e a base de dados relacional normalizada;
    \item Dois scrapers baseados em Playwright para recolha automática de dados do Continente e do Pingo Doce;
    \item Um orchestrator de scrapers baseado em GitHub Actions reusable workflows;
    \item Um worker de notificações que integra a base de dados com o Firebase Cloud Messaging;
    \item Uma aplicação web em React com TypeScript e Vite, com suporte a concorrência otimista nas listas de compras e arquitetura white-label dinâmica injetada em tempo de \textit{build};
    \item Uma aplicação móvel em React Native com Expo, com suporte a leitura de código de barras, localização, notificações push e internacionalização;
    \item Um gateway Nginx como ponto de entrada unificado para as aplicações cliente;
    \item Pipelines de \ac{CI/CD} com GitHub Actions para todos os serviços, com publicação automática no GitHub Container Registry e deployment na plataforma Fly.io;
    \item Mecanismos de concorrência otimista e tratamento centralizado de erros no List Service, com adaptação correspondente nas aplicações cliente;
    \item Pipeline de personalização white label na aplicação web, activável por parâmetros no workflow de \ac{CI/CD} sem alteração ao código-fonte partilhado.
\end{itemize}

\section{Desafios Encontrados}
\label{sec:8.2}

O desenvolvimento do ShopISEL apresentou desafios técnicos relevantes em várias frentes.

\textbf{Recolha e normalização de dados.} A extração de dados dos websites do Continente e do Pingo Doce revelou-se um processo frágil, dado que as estruturas HTML das páginas podem mudar sem aviso. O uso do Playwright permitiu lidar com páginas de renderização dinâmica, mas a manutenção dos scrapers exige atenção contínua às alterações dos websites fontes. A normalização dos dados provenientes de fontes heterogéneas, com formatos e nomenclaturas distintos, exigiu a criação de uma camada de matching (MatchQueue) que identificasse o mesmo produto em diferentes insígnias.

\textbf{Isolamento de autenticação em testes.} A integração com o Keycloak é essencial para a segurança do sistema, mas representa um ponto de dependência externa nos testes. A solução encontrada, com a criação de um \texttt{TestAuthHandler} que simula a autenticação em contexto de teste, permite executar testes de integração completos sem necessidade de uma instância Keycloak externa.

\textbf{Arquitetura multi-repositório.} A gestão de múltiplos repositórios com ciclos de desenvolvimento independentes requer disciplina no versionamento, na comunicação de contratos entre serviços e na coordenação de alterações que afetam múltiplos componentes. A utilização de GitHub Actions e de ficheiros de configuração declarativos (\texttt{render.yaml}) contribuiu para mitigar esta complexidade.

\textbf{Sincronização multi-cliente.} A possibilidade de editar a mesma lista a partir da aplicação web e da aplicação móvel impôs a adoção de concorrência otimista no List Service e a adaptação das aplicações cliente ao novo contrato de versão. Este desafio transversal exigiu alterações coordenadas entre backend e frontend, bem como testes que validassem não só o comportamento isolado da API, mas também a experiência do utilizador perante conflitos de atualização.

\textbf{Localização e contexto de compra.} A funcionalidade de apresentação de preços próximos depende da localização do dispositivo, o que acrescenta valor à decisão do utilizador sem exigir um modelo geográfico pesado no servidor. O sistema pede permissão quando necessário e usa essa informação apenas para ordenar e contextualizar resultados.

\section{Considerações Finais}
\label{sec:8.3}

O projeto ShopISEL demonstra que é possível construir uma plataforma digital funcional e coerente com base em práticas de engenharia de software, mas também evidencia limites concretos do processo seguido. Adoptaram-se decomposição por serviços, contratos explícitos entre componentes, contentorização, \ac{CI/CD}, testes de integração e revisão iterativa das funcionalidades, o que ajudou a manter o sistema evolutivo e verificável.

Ainda assim, o desenvolvimento não correspondeu a um processo A-SDLC completo, tal como enquadrado por Bhati~\cite{bhati-asdlc-2026}. Não houve delegação autónoma de tarefas de ponta a ponta a agentes, nem um sistema formal de governação com registo sistemático de decisões, validação de cada passo por agentes especializados e controlo contínuo do contexto. O uso de LLMs foi sobretudo assistivo, apoiando implementação, refinamento e validação pontual.

Esta distinção é importante: o valor do projeto está menos na adesão total a um paradigma agentic e mais na combinação pragmática entre desenvolvimento orientado por humanos, automação de entrega e apoio pontual por IA. Em comparação com uma abordagem A-SDLC mais madura, o trabalho ainda fica aquém de práticas como orquestração autónoma de tarefas, verificação sistemática multiagente e rastreabilidade completa das decisões assistidas por IA.

Em suma, o ShopISEL representa um primeiro passo sólido para uma plataforma de apoio à decisão de compra que, com as direções de evolução apresentadas no Capítulo~\ref{chap:7}, tem potencial para se tornar uma ferramenta genuinamente útil no quotidiano dos consumidores portugueses.
